\section{The Euler connection and the sharp lattice cokernel}

% BEGIN E85C1
\paragraph{E85C1. An integral first derivative on the crystalline Frobenius columns.}
Use the tame field $K_t$ and coefficient ring $O'$ of Q84T1--2.
Write $W=W(\mathbb F_{p^a})$, $E(U)=U^e+p$, and let $M$ be the
evaluation of the crystalline lattice's prismatic $F$-crystal on
$(W[[U]],E)$. We use the crystalline-lattice equivalence of Bhatt and
Scholze (\emph{Prismatic $F$-crystals and crystalline Galois
representations}, Theorem 1.2). Work in the component bases of this
Breuil--Kisin model. Let $A_\tau$ be its Frobenius matrix, and let
$\mathcal L_\tau$ be its Frobenius-column span over $O'[[u]]$.
We prove that the ordinary connection $\partial$ preserves
$S\mathcal L_\tau$, where $S$ consists of the series $\sum f_nu^n$
such that $\lfloor n/e\rfloor! f_n\in O'$ and these products tend
to zero. This is the precise integral derivative needed below.

First use the same ring with coefficients $W$, denoted $S_W$. Its
weighted coefficients identify its additive group with sequences in
$W$ tending to zero. The ratios
$\lfloor (i+j)/e\rfloor!/(\lfloor i/e\rfloor!\lfloor j/e\rfloor!)$
are integers, so multiplication extends continuously from
polynomials. The ring is $p$-torsion free and $p$-adically complete.
Coefficient Frobenius with $u\mapsto u^p$ lifts Frobenius modulo $p$.
Indeed for $m=\lfloor n/e\rfloor\ge1$ and $k=\lfloor pn/e\rfloor$
we have $v_p(k!)\ge m+v_p(m!)$. Both the image and the $p$-th power
of $u^n/m!$ thus belong to $pS_W$: for the second use
$m-(p-1)v_p(m!)\ge1$ by the digit-sum formula. For $m=0$ the assertion
is ordinary Frobenius congruence. Density of the finite sums proves
the assertion on $S_W$.

Put $D=S_W[t]/(t^2)$ with Frobenius $t\mapsto pu^{p-1}t$.
It too is $p$-torsion free, complete, and Frobenius lifts the
$p$-th power. Hence $(D,(p))$ is a bounded crystalline prism,
with $\delta(x)=(\varphi(x)-x^p)/p$. There are two prism maps from
the Breuil--Kisin prism, acting on $W$ by Frobenius and sending
$U$ to $u^p$ and $u^p+pu^{p-1}t$, respectively. They commute with
Frobenius by expansion modulo $t^2$. Their evaluations on $W[[U]]$
exist since the weighted coefficients of the substituted series and
its first-order term tend to zero. Both images of $E$ belong to
$pD$, since $u^{ep}/p\in S_W$. The two maps agree modulo $p$,
so give the same structural map from $W[U]/(E)=O_{K_t}$ to $D/p$.
Projection $t\mapsto0$ identifies them on $(S_W,(p))$.

The crystal base-change isomorphisms along the two maps give a
transition matrix $1+tB_\tau$ from the second basis to the first.
The constant is the identity by projection to the diagonal. These
isomorphisms commute with the $O'$-action. Indeed, the coefficient
endomorphisms lift by full faithfulness of the crystalline-lattice
equivalence. The splitting of $W\otimes_{\mathbb Z_p}O'$ into the
unramified components identifies the extended scalar ring on each
component with $S[t]/(t^2)$; this follows also by expanding the
weighted null sequences in a finite integral coefficient basis.
In particular every entry of $B_\tau$ belongs to $S$.

Here is the identification with the actual connection. In split
coordinates $\varphi f=f(u^p)$ with a shift of the preceding component.
Put $C_\tau=\varphi(A_{\tau-1})$. Frobenius compatibility of the
two base changes, evaluated over the rational formal power series
where the Frobenius matrices are invertible, reads
\[
 (1+tB_\tau)(C_\tau+tC_\tau')
   =C_\tau(1+pu^{p-1}t\,\varphi(B_{\tau-1})).
\]
There is at most one system of first derivatives solving this
equation. A difference has $u$-adic order at least $p-1$ plus $p$
times the order on the preceding component, since $C_\tau$ and its
inverse are regular at zero. Taking the least order of any nonzero
entry over all components gives a contradiction.

For existence of the rational solution in the canonical coordinates,
let $\Xi_\tau(0)=1$ and
$\Xi_\tau=A_\tau\varphi(\Xi_{\tau-1})A_\tau(0)^{-1}$.
Successive substitution gives these matrices coefficient by
coefficient because $u\mapsto u^p$ raises each positive order.
In the frame $\Xi$ the ordinary connection differentiates coordinates.
Set $W_\tau=\Xi_\tau^{-1}A_\tau$. The identity
$W_\tau C_\tau=A_\tau(0)\varphi(W_{\tau-1})$, upon differentiation,
shows that $W_\tau^{-1}W_\tau'$ solves the equation for $B_\tau$.
Thus the connection in the Frobenius-column basis is precisely
$f\mapsto f'+B_\tau f$. Finally $S$ is closed under differentiation:
the new coefficient $(n+1)f_{n+1}$ has the required weight, since
when $n+1=em$ the factor $n+1$ cancels the factorial ratio $m$,
and otherwise the weights are equal. This proves the asserted
integrality, including the convergence condition.
% END E85C1

% BEGIN E85C2
\paragraph{E85C2. The integral lattice and the terms of degree at least $ep$.}
Let $\mathcal M=S\mathcal L$, $S_{\rm hi}=S\cap u^{ep}L'[[u]]$ and
$\mathcal H=S_{\rm hi}\mathcal L$, component by component. For
$n<ep$, $\lfloor n/e\rfloor!$ is a unit, so truncation of an element
of $S$ before degree $ep$ has integral coefficients. Conversely every
integral series belongs to $S$, since the factorial weights tend to
zero $p$-adically. Therefore, in the Frobenius-column coordinates,
\[
 \mathcal M=\mathcal L+\mathcal H,\qquad
 \mathcal L\cap\mathcal H=u^{ep}\mathcal L.
\]
The second identity follows by shifting the coefficients of an
integral series whose first $ep$ coefficients vanish.
Write $N=u\partial$ for the Euler connection. E85C1 gives
$N\mathcal M\subset u\mathcal M$. Multiplication by $u$ preserves
$S_{\rm hi}$, as does scalar Euler differentiation
$\sum f_nu^n\mapsto\sum nf_nu^n$. In the Frobenius-column formula
the remaining term is $uB_\tau f$, which belongs to $S_{\rm hi}$
coordinatewise if $f$ does. Thus $N\mathcal H\subset\mathcal H$.

An integral Frobenius morphism $G_\tau$ commutes with the connections:
$\Xi_{2,\tau}^{-1}G_\tau\Xi_{1,\tau}$ is constant, since Frobenius
equivariance conjugates the same expression on the preceding
component at $u^p$ by the constant Frobenius matrices. Subtract the
constant terms and iterate in positive degree to prove this assertion.
Its action in Frobenius-column coordinates is
$A_{2,\tau}^{-1}G_\tau A_{1,\tau}=\varphi(G_{\tau-1})$, with integral
coefficients. Hence it takes $\mathcal H_1$ into $\mathcal H_2$.
For a tame descent automorphism the same horizontal-frame argument
applies with the rescaling $u\mapsto\zeta_\tau u$ on its source.
Scalar Euler differentiation commutes with this rescaling, and
the scalar and matrix actions preserve the high-degree ideal and
the integral Frobenius-column lattice. Thus $N$ preserves every
tame isotypic part.
% END E85C2

% BEGIN E85C3
\paragraph{E85C3. Degree raising on the actual tame character spaces.}
Fix the componentwise tame index $c\in\mathbb Z/e$. Project by
averaging over $\mu_e$ against the inverse character on each
component. All coefficients are integral since $e$ is a unit and
the character values are roots of unity in $O'$. The projectors
sum to the identity by character orthogonality. They preserve
$u^i\mathcal L$, $\mathcal H$ and commute with $G$ and $N$.
Put $V_i=(u^i\mathcal L+\mathcal H)_c$. Every element of $V_i$ is
$u^iy+z$, with $y\in\mathcal L$, $z\in\mathcal H$, and with each
summand in the $c$ part: project a decomposition to obtain one.
By Leibniz and E85C2,
\[
 (N-i)(u^iy)=u^iNy\in u^{i+1}\mathcal M
       \subset u^{i+1}\mathcal L+\mathcal H .
\]
The high-degree part is stable, so $(N-i)V_i\subset V_{i+1}$.

If $c-i\notin A$ there is a stronger inclusion $V_i\subset V_{i+1}$.
In fact the nonzero constant term of $y$ in any component would
transform under inertia as $\zeta_\tau^{c-i}$, by cancellation of
$u^i$ in the eigencharacter identity and evaluation at zero.
The tame type of the model restricts its characters on that
fibre to $A$, contrary to the hypothesis. Consequently $y(0)=0$.
Writing $y=A_\tau d$ in each Frobenius-column basis gives $d(0)=0$
because $A_\tau(0)$ is invertible over $L'$. Division of the
integral coefficients $d$ by $u$ now gives $y\in u\mathcal L$.

In the notation of Q84T5 traverse the integers $0\le i<eJ$.
At a represented value $i=(c-s)\bmod e+je$, $s\in A$, apply $N-i$.
At each other value use the stronger inclusion and apply the identity.
This gives exactly $P_1(N)\mathcal L_c\subset
(u^{eJ}\mathcal L+\mathcal H)_c$. Starting at $i=eJ$ and ending at
$2eJ-1$ proves $P_2(N)(u^{eJ}\mathcal L)_c\subset
(u^{2eJ}\mathcal L+\mathcal H)_c$. In particular multiplicities of
the tame characters introduce no additional factors.
% END E85C3

% BEGIN E85C4
\paragraph{E85C4. The cokernel on the Frobenius columns and invariant lattices.}
For a realized morphism of the type, height and finite-length
cokernel in Q84T2--3, put $v=u^e$, $B=w+1$ and $J=H+b+B$ as in
Q84T5. Its error module
\[
 \mathcal T=v^JG\mathcal L_1+p^B\mathcal L_2+\mathcal H_2
\]
is $N$-stable. On the first term the scalar derivative contributes
$eJ$ times the same term. The other contribution belongs to
$v^JG(\mathcal L_1+\mathcal H_1)$ by horizontality and E85C2.
On the second term the derivative belongs to
$p^B(\mathcal L_2+\mathcal H_2)$. Both errors belong to
$\mathcal H_2$, which is stable. Thus the two degree raising
statements, the high-degree intersection and the error stability
required by Q84T5 hold on this very morphism. F84R2 supplies its
coarse congruence. Apply the integral interpolation and cancellation
in Q84T5 and the coefficient extraction and integral averaging in
F84R3. We obtain
\[
 p^w\mathcal L_2\subset G\mathcal L_1,\qquad
 p^w M_2\subset GM_1,\qquad p^wY'_2\subset GY'_1 .
\]
The same bound holds on $E^{-b}M$, $M^\Delta$ and
$Y=(E^{-b}M)^\Delta$. This supplies the sharp cokernel estimate in
Q84T4, whose separate uses of period representatives and correction
of the monodromy term require their own estimates.
% END E85C4
